\clearpage
\begin{trapbox}
  \textbf{\textcolor{CTrap}{易错 1：$\Delta(y)\ge0$ 解出了闭区间，就直接把两端都写进去。}}
\begin{center}\begin{tikzpicture}[font=\small,>=stealth,line join=round]
\draw[gray!45,rounded corners=3pt] (0,0) rectangle (14.2,9.4);
\node at (7.1,8.8) {函数无限接近下边界，但从不落到下边界上};

\begin{scope}[shift={(4.1,1.8)},x=.9cm,y=1.65cm]
\draw[->] (-3.5,0)--(3.6,0) node[right]{$x$};\draw[->] (0,-.1)--(0,3.8) node[above]{$y$};\draw[dashed,magenta] (-3.3,2)--(3.3,2);
\draw[teal!80!black,thick,domain=-3.2:3.2,samples=100] plot(\x,{2+1/(\x*\x+1)});\fill[blue] (0,3) circle(2pt);\node[above right] at (0,3){$(0,3)$};\node[above left] at (0,2){$2$};
\end{scope}
\node[align=left,anchor=west] at (8.3,5.1){$y=2+\dfrac1{x^2+1}$\\[10pt] $0<\dfrac1{x^2+1}\le1$\\[10pt] 值域 $(2,3]$\\[12pt] $y=2$ 时变成 $1=0$};
\node[text width=13.2cm,align=center] at (7.1,.65) {即使候选判别式在 $y=2$ 处为零，降次后也可能无解。};
\end{tikzpicture}
\end{center}


  \medskip
  \textbf{错在哪里：}当 $a-dy=0$ 时，方程已经不是二次方程；这个 $y$ 即使满足写出来的 $\Delta(y)\ge0$，也不一定取到。
  例如 $y=\dfrac{2x^2+3}{x^2+1}$，照搬判别式会得到 $[2,3]$，但 $y=2$ 时要求 $1=0$，不可能成立。
  \par\textbf{正确做法：}先用判别式求候选范围，再另查 $y=\dfrac ad$；正确答案是 $(2,3]$。

  \bigskip
\end{trapbox}
\clearpage
\begin{trapbox}
\textbf{\textcolor{CTrap}{易错 2：分母可能等于 $0$，还认为“方程有解就一定行”。}}
\begin{center}\begin{tikzpicture}[font=\small,>=stealth,line join=round]
\draw[gray!40,rounded corners=3pt] (0,0) rectangle (14.2,7.8);
\node at (7.1,7.2) {去分母后得到的根，可能恰好落在原函数的孔点};
\begin{scope}[shift={(4.4,2.0)},x=1cm,y=2cm]
\draw[->,gray] (-3.4,0)--(3.4,0) node[right]{$x$};\draw[->,gray] (0,-.2)--(0,2.2) node[above]{$y$};
\draw[blue,very thick] (-3.1,1)--(3.1,1);
\draw[orange!85!black,thick,fill=white] (-1,1) circle(3pt);\draw[orange!85!black,thick,fill=white] (1,1) circle(3pt);
\node[above left=7pt] at (-1,1){$(-1,1)$};\node[above right=7pt] at (1,1){$(1,1)$};
\node[below left=4pt] at (0,0){$O$};
\end{scope}
\node[align=left,anchor=west] at (8.4,4.2){$y=\dfrac{x^2-1}{x^2-1}$\\[12pt]$x\ne\pm1$\\[12pt]两处孔点不能当作解\\[8pt]真正值域仍为 $\{1\}$};
\node[text width=13cm,align=center] at (7.1,.6) {空心点对应分母为零的位置；它们不能证明其他 $y$ 可以取到。};
\end{tikzpicture}
\end{center}


  \medskip
  \textbf{错在哪里：}分式分母不能为 $0$，乘过去以后有时会出现原来不允许的 $x$。
  例如
  \[
    y=\frac{x^2-1}{x^2-1},\qquad x\ne\pm1.
  \]
  真正的函数值始终是 $1$。但若只写 $(1-y)(x^2-1)=0$，误把 $x=\pm1$ 当成解，就会错误地说任意 $y$ 都能取到。
  \par\textbf{正确做法：}先查分母。若分母始终不为 $0$，可以放心用本结论；否则每个找到的 $x$ 都要确认没有让原分母等于 $0$。

  \bigskip
\end{trapbox}
\clearpage
\begin{trapbox}
\textbf{\textcolor{CTrap}{易错 3：看见 $x^2$ 就套判别式。}}
\begin{center}\begin{tikzpicture}[font=\small,>=stealth,line join=round]
\draw[gray!40,rounded corners=3pt] (0,0) rectangle (14.2,13.5);
\node at (7.1,12.9) {二次项消失之后：改看剩余式子 $H(x)$ 有没有零点};
\begin{scope}[shift={(3.6,10.3)},x=1cm,y=.8cm]
\draw[->,gray] (-2.2,0)--(2.3,0) node[right]{$x$};
\draw[->,gray] (0,-1.5)--(0,1.65) node[above]{$H$};
\draw[blue,very thick,domain=-1.1:.1] plot(\x,{-2*\x-1});\fill[teal!80!black] (-.5,0) circle(2.2pt);
\end{scope}
\node[align=left,anchor=west] at (7.3,10.45) {$y=1$：$-2x-1=0$\\[8pt]与横轴相交，$x=-1/2$ 可用};
\begin{scope}[shift={(3.6,6.6)},x=1cm,y=.8cm]
\draw[->,gray] (-2.2,0)--(2.3,0) node[right]{$x$};
\draw[->,gray] (0,-1.5)--(0,1.65) node[above]{$H$};
\draw[blue,very thick] (-2,1)--(2,1);
\end{scope}
\node[align=left,anchor=west] at (7.3,6.75) {$y=2$：$1=0$\\[8pt]始终离开横轴，无解};
\begin{scope}[shift={(3.6,2.9)},x=1cm,y=.8cm]
\draw[->,gray] (-2.2,0)--(2.3,0) node[right]{$x$};
\draw[->,gray] (0,-1.5)--(0,1.65) node[above]{$H$};
\draw[blue,very thick] (-2,0)--(2,0);
\end{scope}
\node[align=left,anchor=west] at (7.3,3.05) {$y=2$：$0=0$\\[8pt]整条横轴都是零点，任意 $x$ 可用};

\node[text width=13cm,align=center] at (7.1,.6) {从上到下对应标准题、缺端点题、常函数题；仍须确认原分母不为零。};
\end{tikzpicture}
\end{center}


  \medskip
  \textbf{错在哪里：}判断 $\Delta\ge0$ 只适用于二次项系数\textbf{不为 $0$}的一元二次方程。
  \par\textbf{正确做法：}先看 $(a-dy)x^2$ 前面的 $a-dy$。如果它等于 $0$，就改成解一次方程，或者检查有没有出现 $0=0$ 与 $1=0$ 这种情况。

  \bigskip
\end{trapbox}
\clearpage
\begin{trapbox}
\textbf{\textcolor{CTrap}{易错 4：把“越来越接近”当成“能够等于”。}}
\begin{center}\begin{tikzpicture}[font=\small,>=stealth,line join=round]
\draw[gray!45,rounded corners=3pt] (0,0) rectangle (14.2,9.4);
\node at (7.1,8.8) {函数无限接近下边界，但从不落到下边界上};

\begin{scope}[shift={(4.1,1.8)},x=.9cm,y=1.65cm]
\draw[->] (-3.5,0)--(3.6,0) node[right]{$x$};\draw[->] (0,-.1)--(0,3.8) node[above]{$y$};\draw[dashed,magenta] (-3.3,2)--(3.3,2);
\draw[teal!80!black,thick,domain=-3.2:3.2,samples=100] plot(\x,{2+1/(\x*\x+1)});\fill[blue] (0,3) circle(2pt);\node[above right] at (0,3){$(0,3)$};\node[above left] at (0,2){$2$};
\end{scope}
\node[align=left,anchor=west] at (8.3,5.1){$y=2+\dfrac1{x^2+1}$\\[10pt] $0<\dfrac1{x^2+1}\le1$\\[10pt] 值域 $(2,3]$\\[12pt] $y=2$ 时变成 $1=0$};
\node[text width=13.2cm,align=center] at (7.1,.65) {即使候选判别式在 $y=2$ 处为零，降次后也可能无解。};
\end{tikzpicture}
\end{center}


  \medskip
  对于 $d\ne0$ 的二次分式，$|x|$ 越来越大时，函数值越来越接近 $\dfrac ad$，但这个数未必能取到。
  \par\textbf{快速检查：}直接看
  \[
    \frac{ax^2+bx+c}{dx^2+ex+f}-\frac ad
    =\frac{(bd-ae)x+(cd-af)}{d(dx^2+ex+f)}.
  \]
  分母始终不为 $0$ 时，只需问分子能不能等于 $0$；这正好回到方法一中的方程 (3)。

  \bigskip
\end{trapbox}
\clearpage
\begin{trapbox}
\textbf{\textcolor{CTrap}{易错 5：展开 $\Delta(y)$ 时丢掉负号。}}
\begin{center}\begin{tikzpicture}[font=\small,>=stealth,line join=round]
\draw[gray!40,rounded corners=3pt] (0,0) rectangle (14.2,9.2);
\node at (7.1,8.6) {漏掉负号会把开口翻转，候选范围也会选反};
\pgfmathsetmacro{\LL}{(3-sqrt(5))/2}
\pgfmathsetmacro{\UU}{(3+sqrt(5))/2}
\pgfmathsetmacro{\rr}{(sqrt(5)-1)/2}
\pgfmathsetmacro{\ss}{-(sqrt(5)+1)/2}
\begin{scope}[shift={(2.5,4.2)},x=1.65cm,y=.40cm]
\draw[->,gray] (-.2,0)--(3.1,0) node[above=4pt]{$y$};\draw[->,gray] (0,-5.6)--(0,5.8) node[above]{$\Delta$};
\draw[blue,very thick,domain=.15:2.85,samples=120] plot(\x,{4*(-\x*\x+3*\x-1)});
\draw[orange!85!black,dashed,thick,domain=.15:2.85,samples=120] plot(\x,{-4*(-\x*\x+3*\x-1)});
\draw[dotted,gray] (\LL,0)--(\LL,-5.7);\draw[dotted,gray] (\UU,0)--(\UU,-5.7);\node[below=6pt] at (\LL,-5.7){$L$};\node[below=6pt] at (\UU,-5.7){$U$};
\end{scope}
\node[align=left,anchor=west] at (8.4,5.2){\textcolor{blue}{正确：开口向下}\\[7pt]$4(-y^2+3y-1)$\\[12pt]\textcolor{orange!85!black}{示例错误：整个式子反号}\\[7pt]$-4(-y^2+3y-1)$\\[12pt]“中间非负”会误成“两侧非负”};
\node[text width=13cm,align=center] at (7.1,.6) {本例分母判别式为 $-4<0$；正确 $\Delta(y)$ 的 $y^2$ 系数也应是 $-4$。};
\end{tikzpicture}
\end{center}


  \medskip
  建议保留中间一步：
  \[
    \begin{aligned}
      &(b-ey)^2-4(a-dy)(c-fy)\\
      &=(b^2-2bey+e^2y^2)\\
      &\quad-4(ac-afy-cdy+dfy^2).
    \end{aligned}
  \]
  最终 $y^2$ 的系数应该是 $e^2-4df$。若题目已知 $e^2-4df<0$，算出来的系数却是正数，就该重新检查符号。


  % ==== 新增：几何方法需要的防错检查 ====
  \bigskip
\end{trapbox}
\clearpage
\begin{trapbox}
\textbf{\textcolor{CTrap}{易错 6：把“只有一个交点”都当成相切。}}
\begin{center}\begin{tikzpicture}[font=\small,>=stealth,line join=round]
\draw[gray!40,rounded corners=3pt] (0,0) rectangle (14.2,9.8);
\node at (7.1,9.200000000000001) {只有一个有限交点，也可能是穿过曲线而非相切};
\begin{scope}[shift={(1.6,1.8)},x=.83cm,y=.83cm]
\draw[->,gray] (0,0)--(6.5,0) node[right]{$X$};\draw[->,gray] (0,0)--(0,6.7) node[above]{$Y$};
\draw[blue,very thick,variable=\t,domain=-2.3:1.1,samples=160] plot({\t*\t+2*\t+2},{\t*\t+1});
\draw[orange!85!black,thick] (0,0)--(5.8,5.8);
\fill[orange!85!black] (1.25,1.25) circle(2.2pt) node[below=12pt]{$P$};
\node[below left=4pt] at (0,0){$O$};
\end{scope}
\node[align=left,anchor=west] at (8.1,5.2){橙线：$Y=X$\\[12pt]$-2x-1=0$\\[8pt]只有 $x=-1/2$\\[12pt]$P=(5/4,5/4)$\\[12pt]直线穿过曲线，并非相切};
\node[text width=13cm,align=center] at (7.1,.6) {这里关于原自变量 $x$ 的方程已是一次方程，不能用二次方程的重根判断相切。};
\end{tikzpicture}
\end{center}


  \medskip
  对 $y=\dfrac{x^2+1}{x^2+2x+2}$，令 $Y=kX$ 后得到
  $(1-k)x^2-2kx+(1-2k)=0$。
  当 $k=1$ 时，它只剩下 $-2x-1=0$，虽然恰好找到一个有限的 $x$，
  却\textbf{不是}由二次方程判别式等于 $0$ 得出的切点。
  \par\textbf{正确做法：}只有关于 $x$ 的二次项确实存在时，才用 $\Delta=0$ 表示边界相切；二次项消失要单查。

  \bigskip
\end{trapbox}
\clearpage
\begin{trapbox}
\textbf{\textcolor{CTrap}{易错 7：把整个单位圆都当成原来的 $x$ 可以走到。}}
\begin{center}\begin{tikzpicture}[font=\small,>=stealth,line join=round]
\draw[gray!40,rounded corners=3pt] (0,0) rectangle (14.2,9);
\node at (7.1,8.4) {圆上恰好少了最右点：值域的左端点因此不取};
\begin{scope}[shift={(3.8,4.3)},x=2.1cm,y=2.1cm]
\draw[->,gray] (-1.3,0)--(1.4,0) node[right]{$u$};\draw[->,gray] (0,-1.3)--(0,1.4) node[above]{$v$};
\draw[blue,very thick] (0,0) circle(1);
\fill[teal!80!black] (-1,0) circle(2.2pt) node[above left=6pt]{$P$};
\draw[orange!85!black,thick,fill=white] (1,0) circle(3pt) node[above right=6pt]{$E$};
\end{scope}
\node[align=left,anchor=west] at (8.1,5.1){$y=\tfrac52-\tfrac12u$\\[12pt]\textcolor{teal!80!black}{$P=(-1,0)$：$y=3$，取到}\\[12pt]\textcolor{orange!85!black}{$E=(1,0)$：$y=2$，取不到}\\[12pt]$-1\le u<1$\\[8pt]$2<y\le3$};
\node[text width=13cm,align=center] at (7.1,.6) {$P$ 对应 $x=0$；$E$ 只能无限接近，不能由有限实数 $x$ 得到。};
\end{tikzpicture}
\end{center}


  \medskip
  对分母 $D=x^2+2x+2$ 所构造的圆，
  \[
    u=\frac{D-2}{D},\qquad v=\frac{2x+2}{D},
    \qquad 1-u=\frac2D>0.
  \]
  所以 $u=1$、$v=0$ 这个圆上的点\textbf{永远不可能由有限的 $x$ 得到}。
  \par\textbf{正确做法：}直线与整个圆相交，只给出了候选范围；
  若边界碰到不能到达的点，该值要去掉。例 1 的 $y=2$ 就是这种情况。

  \bigskip
\end{trapbox}
\clearpage
\begin{trapbox}
\textbf{\textcolor{CTrap}{易错 8：只要看见二次式相除，就直接断言轨迹一定是抛物线或一定能凑成圆。}}
\begin{center}\begin{tikzpicture}[font=\small,>=stealth,line join=round]
\draw[gray!40,rounded corners=3pt] (0,0) rectangle (14.2,9.4);
\node at (7.1,8.8) {分子分母成比例：点在移动，连线斜率却保持不变};
\begin{scope}[shift={(3.3,1.8)},x=.75cm,y=.75cm]
\draw[->,gray] (0,0)--(3.4,0) node[right]{$X$};\draw[->,gray] (0,0)--(0,7.5) node[above]{$Y$};
\draw[blue,very thick,->] (1,2)--(3.3,6.6);
\fill[blue] (1,2) circle(2.2pt) node[above left=5pt]{$(1,2)$};
\node[below left=4pt] at (0,0){$O$};
\end{scope}
\node[align=left,anchor=west] at (8,5.1){$X=x^2+1\ge1$\\[9pt]$Y=2x^2+2=2X$\\[12pt]动点沿射线运动\\[8pt]原点连线斜率始终是 $2$\\[12pt]轨迹不是抛物线\\[8pt]函数值也不随点移动而变};
\node[text width=13cm,align=center] at (7.1,.6) {横纵坐标等比例；当 $d\ne0$ 时横坐标会变化，轨迹不会固定成一个点。};
\end{tikzpicture}
\end{center}


  \medskip
  斜率法可以从分母、分子分别作为 $X,Y$ 自然出发，但动点有时可能沿一条直线的一部分（例如射线）运动；本条 $d\ne0$，所以横坐标会变化，不会固定成一个点；
  圆则需要额外找到真正成立的勾股恒等式。
  \par\textbf{正确做法：}先写出 $X,Y$ 并检查图形，再决定怎样求斜率；
  若想用圆，先从原分母用平方差\textbf{推导出}圆方程，不能倒过来生造两个分式。
\end{trapbox}
